\chapter*{学士論文概要}
\addcontentsline{toc}{chapter}{学士論文概要}
%近年， SNSの利用者はますます増加しておりSNSは社会のインフラとして欠かせなくなっている．SNSは主にコミュニケーションや情報収集に利用されている．また，SNSには速報性があり，ユーザはタイムリーに社会的な関心事や自分の興味に基づいた情報を入手できる．しかし自分が興味ある内容だからといって，必ずしも今必要な情報だとは限らない．その代表例が「ネタバレ」である．例えば，スポーツ中継をリアルタイムでの視聴ができないため，仕方なく録画予約をして後で視聴する人もいる．そのような視聴者が，視聴前にそのスポーツの結果に関連するネタバレツイートを見てしまうと，試合の面白さを低下させる恐れがある．
近年， SNSの利用者はますます増加しておりSNSは社会のインフラとして欠かせなくなっている．SNSは主にコミュニケーションや情報収集に利用されている．SNSの代表例としてTwitterがあり，Twitterは他のアカウントをフォローすることで，そのアカウントのツイート（投稿）を見ることができる．Twitterの用途は人それぞれであり，例えば，あるユーザは親しいユーザをフォローし，そのユーザとコミュニケーションを行ったり，またあるユーザは自分の興味がある分野に関する情報を積極的に投稿するアカウントをフォローし，情報収集を行っている．さらにTwitterには速報性があり，ユーザはタイムリーに社会的な関心事や自分の興味に基づいた情報を入手できる．しかし自分が興味ある内容だからといって，必ずしも今必要な情報だとは限らない．その代表例が「ネタバレ」である．例えば，スポーツ中継をリアルタイムで視聴できないため，録画をして後で視聴する人もいる．そのような視聴者が，視聴前にそのスポーツの結果に関連するネタバレツイートを見てしまうと，試合の面白さを低下させる恐れがある．

そのため本研究では，Twitter上のネタバレ，特にサッカーの試合に関するネタバレを自動で検出し，閲覧を防止するモデルを構築する．さらに，検出するネタバレをスコアに関係するネタバレとスコアに関係しないネタバレに分類する．これにより試合の重要な情報を抽出することが可能になり，本モデルは，ネタバレ防止の他に，試合速報やハイライトなど，試合情報の自動要約に応用されることが見込まれる．


評価実験では，Twitterからツイートを取得し，自作辞書によって形態素解析器MeCabに認識されない用語を別の用語に置換した．また，齊藤ら \cite{齊藤}が用いた「興奮度」をテキスト内の「ー」，「!」，「あ」，「ぁ」という文字の出現数の合計とし，特徴量として追加した．そしてナイーブベイズ，SVM，ランダムフォレストの3種類の機械学習を用いて分類を行った結果，ランダムフォレストのマクロ平均適合率は0.800，マクロ平均再現率は0.764，マクロ平均F1値は0.778となり，ランダムフォレストが本研究における最適なモデルとなった．また，特徴量の重要度を算出しランク付けしたところ，自作辞書によって置換された用語や興奮度がランキングの上位に存在した．そのため，自作辞書と興奮度はモデルの性能向上に貢献したと考えられる．%今後の展望として，データ数を増やし，モデルが様々な試合展開に対応できるようにすることや，本研究で分類した様々な種類のネタバレ情報をもとに，試合速報や試合のハイライトなどを自動で作成することなどが挙げられる．